\chapter*{Preparing Small Molecule Library for Validation and Virtual Screening}
One of the most important, albeit neglected, steps in molecular docking is validation. If the receptor was co-crystallised with a ligand in the receptor site of interest, the protocol must be validated by re-docking it and checking if the experimental binding conformation can be reproduced. However, when there are no bound ligands reported, we advise gathering from the literature a few molecules targeting the receptor in question with experimental binding affinity values reported. After docking these molecules, a rank ordering is carried out to understand if the docking parameters can reproduce the experimental binding affinities in the correct order. We strongly recommend that trying to reproduce the binding affinity values as determined by experimental methods is extremely difficult to inherent limitations in the scoring functions. These limitations are outlined briefly in the previous section and we also direct the readers towards a comprehensive article for further explanation\cite{bib54} A large collection of data can be found in the \href{https://www.bindingdb.org/rwd/bind/index.jsp}{BindingDB} or  \href{https://www.ebi.ac.uk/chembl/}{ChemBL}. However, it must be noted that these data may contain results from different experimental protocols making a direct comparison difficult.

Another important yet neglected question that is not often raised is on the number poses that must be generated to ensure better predictions. The number of poses generated in molecular docking varies based on the docking software and the objective of the study. Generally, protocols aim for 10–20 poses per ligand (in our group the standard practice is to generate 10 poses minimum)\cite{bib10} to explore binding conformations effectively. While working with completely new scaffolds, it becomes difficult to establish which docking pose would be more physiologically relevant. In such case, we recommend evaluating at least 2-3 diverse poses and them using MD simulations followed by MM-GB/PB-SA methods for estimating the binding affinity.
